\section{Results}

\begin{table*}[t]
  \caption{Fold-10 Accuracy and Macro F1 Under Controlled Noise}
  \label{tab:complete-test-results}
  \centering
  \begin{tabular}{llrrrrrrrr}
    \toprule
    & & \multicolumn{4}{c}{Accuracy} & \multicolumn{4}{c}{Macro F1} \\
    \cmidrule(lr){3-6}\cmidrule(lr){7-10}
    Architecture & Training & Clean & 20 dB & 10 dB & 0 dB
      & Clean & 20 dB & 10 dB & 0 dB \\
    \midrule
    CNN & Baseline
      & 0.7121 & 0.6703 & 0.6057 & 0.4504
      & 0.7323 & 0.6624 & 0.5857 & 0.4230 \\
    CNN & Augmented
      & 0.7288 & 0.7168 & 0.6989 & 0.5496
      & 0.7469 & 0.7270 & 0.7071 & 0.5481 \\
    CRNN & Baseline
      & 0.7838 & 0.7491 & 0.7085 & 0.4994
      & \textbf{0.8037} & 0.7632 & 0.7134 & 0.4952 \\
    CRNN & Augmented
      & 0.7419 & 0.7491 & 0.7384 & 0.6141
      & 0.7574 & 0.7634 & \textbf{0.7535} & 0.6186 \\
    ResNet18 & Baseline
      & 0.7395 & 0.6977 & 0.5938 & 0.4421
      & 0.7617 & 0.7050 & 0.5774 & 0.4249 \\
    ResNet18 & Augmented
      & 0.7575 & 0.7575 & 0.7264 & 0.6081
      & 0.7781 & \textbf{0.7783} & 0.7432 & \textbf{0.6384} \\
    \bottomrule
  \end{tabular}
\end{table*}

\begin{table*}[t]
  \caption{Macro-F1 Robustness Summary}
  \label{tab:robustness-summary}
  \centering
  \begin{tabular}{llrrrr}
    \toprule
    Architecture & Training & Clean F1 & 0 dB F1 & Normalized SNR AUC & AUC rank \\
    \midrule
    CNN & Baseline & 0.7323 & 0.4230 & 0.5642 & 6 \\
    CNN & Augmented & 0.7469 & 0.5481 & 0.6724 & 3 \\
    CRNN & Baseline & \textbf{0.8037} & 0.4952 & 0.6713 & 4 \\
    CRNN & Augmented & 0.7574 & 0.6186 & 0.7222 & 2 \\
    ResNet18 & Baseline & 0.7617 & 0.4249 & 0.5712 & 5 \\
    ResNet18 & Augmented & 0.7781 & \textbf{0.6384} & \textbf{0.7258} & 1 \\
    \bottomrule
  \end{tabular}
\end{table*}

\begin{table*}[t]
  \caption{Augmented-Minus-Baseline Effects}
  \label{tab:augmentation-effects}
  \centering
  \begin{tabular}{lrrrrr}
    \toprule
    Architecture & Clean accuracy & Clean macro F1
      & 0 dB accuracy & 0 dB macro F1 & Macro-F1 AUC \\
    \midrule
    CNN & +0.0167 & +0.0146 & +0.0992 & +0.1252 & +0.1082 \\
    CRNN & -0.0418 & -0.0464 & +0.1147 & +0.1234 & +0.0509 \\
    ResNet18 & +0.0179 & +0.0163 & +0.1661 & +0.2135 & +0.1546 \\
    \bottomrule
  \end{tabular}
\end{table*}


\ProjectFigure
  {figures/phase_03/fig_p03_03_macro_f1_robustness_curves.png}
  {\columnwidth}
  {Macro-F1 robustness curves}
  {Macro-F1 performance under clean and controlled background-noise
  conditions. Solid lines denote augmented training and dashed lines denote
  baseline training.}
  {fig:macro-f1-robustness}

\ProjectFigure
  {figures/phase_03/fig_p03_05_macro_f1_augmentation_effects.png}
  {\columnwidth}
  {Augmentation-effect comparison}
  {Architecture-matched macro-F1 difference between augmented and baseline
  training at each evaluation condition.}
  {fig:augmentation-effects}

\ProjectFigure
  {figures/phase_03/fig_p03_06_macro_f1_robustness_auc.png}
  {\columnwidth}
  {Robustness-AUC comparison}
  {Normalized macro-F1 SNR area under the noisy-condition curve. Higher values
  indicate stronger performance across 0, 10, and 20 dB.}
  {fig:robustness-auc}

CRNN baseline obtained the highest clean macro F1 (0.8037). ResNet18 augmented
obtained the highest normalized macro-F1 SNR AUC (0.7258) and the highest 0 dB
macro F1 (0.6384). The three augmented configurations occupied the top three
robustness-AUC ranks.

\ProjectFigure
  {figures/phase_03/fig_p03_07_zero_db_per_class_effects.png}
  {\columnwidth}
  {Zero-dB class-level effects}
  {Per-class F1 change from augmentation at 0 dB for each architecture.}
  {fig:zero-db-class-effects}

\ProjectFigure
  {figures/phase_03/fig_p03_08_confusion_matrix_comparison.png}
  {\columnwidth}
  {Normalized clean and 0 dB confusion-matrix comparison -- pending}
  {Normalized confusion matrices for the selected baseline and augmented
  configurations under clean and 0 dB conditions.}
  {fig:confusion-comparison}

\DraftingNote{Expand this section by describing the tables and figures without
repeating every cell. Do not use statistically significant language until
confidence intervals or tests exist.}
